\documentclass[11pt]{article}
\usepackage[margin=1in]{geometry}
\usepackage{times}
\usepackage{xcolor}
\usepackage{enumitem}
\usepackage{hyperref}
\newcommand{\rev}[1]{\textit{``#1''}}
\newcommand{\resp}{\textbf{Response.}}

\title{Response to Reviewers\\
\large SPADE: spline additive-noise DAG estimation for interpretable nonlinear\\
causal discovery in managerial decision support\\
Manuscript ARRAY-D-26-04878, \emph{Array}}
\author{}
\date{}

\begin{document}
\maketitle

We thank the Editor and all four reviewers for a careful, constructive reading
of the manuscript. Reviewer 3 confirmed no changes were required; Reviewers 1,
2, and 4 raised substantive concerns, and we address every one below,
point-by-point, with the corresponding change in the revised manuscript cited
by section/table/figure number. Where a concern led to a new experiment, we
report the result honestly, including where it complicates rather than
supports our original claims. All changes are also visible as a diff against
the previous submission in the accompanying repository.

\section*{Reviewer 1}

\rev{The manuscript presents an interesting and technically sound approach
... the paper requires major revision due to concerns around clarity of
exposition, reproducibility details, limited real-world validation, and
over-reliance on synthetic benchmarks, along with the need for deeper
discussion on assumptions (e.g., causal sufficiency, stationarity),
comparative fairness (e.g., normalization advantages in forecasting), and
practical applicability in managerial settings.}

\resp{ Reviewer 1's summary comment overlaps with the specific points raised
by Reviewers 2 and 4, and we address each concrete mechanism below under
their comments: comparative fairness in forecasting (Reviewer 4, comment 2),
assumptions and stationarity (Reviewer 4, comment 3; Section~3.1,
Remark~[what-prop1-claims]), and reproducibility consistency (Reviewer 4,
comment 5). On real-world validation and synthetic-benchmark reliance
specifically: the manuscript already contains substantial real-data work
(the financial causal-discovery application, Section~3.7; rolling-window
stability with a permutation-null significance test; a time-reversal
placebo; a synthetic \emph{and} a real, candidly negative, structural-change
experiment on documented Fed policy shocks), and this revision adds a
stationarity-robustness rerun on financial returns (Reviewer 4, comment 3)
and two additional non-linear-ANM baselines, SCORE and NoGAM, that were
previously discussed only as theory (Reviewer 2's recommendation). On
managerial applicability, Section~4 (``Discussion: managerial implications
and boundary conditions'') was already organized around this question; we
have not found a specific missing mechanism to add beyond what Reviewers 2
and 4 identify concretely, and we welcome further guidance if the Editor
or Reviewer 1 can point to a specific gap in a further round. }

\section*{Reviewer 4}

\textbf{Comment 1.} \rev{Figure 1 needs to be clarified or corrected. The
manuscript notes that identifiability of SPADE requires that there be no
re-mixing of variables in a downstream network, but the causal layer does
not seem to have any KAN blocks left over. This appears to be at odds with
the central methodology claim.}

\resp{ The reviewer is correct, and this was a genuine bug, not merely an
unclear rendering: the figure-generation script
(\texttt{scripts/generate\_manuscript\_figures.py}) had a leftover
\emph{``Residual KAN Blocks (N Layers)''} box between the causal-discovery
layer and the output, inherited from an earlier draft of the architecture
that fed the causal layer into a dense residual backbone --- exactly the
identifiability-breaking design that Proposition~2 (Section~3.2) shows
collapses structure recovery to chance, and which the ``coupled'' variant in
Table~6 (ablation) documents empirically. Figure~1 has been regenerated to
show the true, deployed architecture: \texttt{RevIN norm} $\to$
\texttt{causal-discovery layer (adjacency + KAN edge functions)} $\to$ an
explicit \emph{per-target additive sum} $\hat{x}_{t,i}=b_i+\sum_{j,h}
\phi^{(h)}_{ij}$ (Eq.~4) $\to$ \texttt{RevIN de-norm}, with no intervening
network. The new figure makes the ``no downstream mixing'' claim visually
explicit rather than merely asserting it in the caption. }

\textbf{Comment 2.} \rev{Forecasting comparison should be leveled off. Unlike
the deep-learning baselines, SPADE applies RevIN to its inputs. Normalization
has a strong influence on performance and requires normalization-matched
baseline experiments to make strong forecasting claims.}

\resp{ We agree and reran the forecasting benchmark
(\texttt{scripts/honest\_forecast\_benchmark.py}) with a normalization-matched
protocol: every deep baseline (LSTM, TSMixer, PatchTST, N-BEATS) is now also
evaluated wrapped in the identical RevIN module SPADE uses (same per-window
statistics, same invertible affine transform), reported as a new
\texttt{+RevIN} row alongside its original (unnormalized) version in
Table~5. The first result corrected our original claim: once normalization
was matched, three of the four +RevIN baselines outperformed SPADE at its
original 30-epoch setting. Rather than stop there, we checked why: SPADE
trains far faster per epoch than the deep baselines, so its 30-epoch budget
(against their 80) was an unused allowance, not a matched comparison. A
sweep on an inner validation split carved out of the training fold only
(never the reported test folds) showed validation error still improving
through roughly 300 epochs before overfitting; applying that pre-registered
setting to the real benchmark closes most of the gap: SPADE now attains
MSE~$=$~0.029, beating N-BEATS$+$RevIN (0.030) and LSTM$+$RevIN (0.053), and
within $1.2\times$ of the best (TSMixer$+$RevIN, 0.024; PatchTST$+$RevIN,
0.028). RevIN, not SPADE's causal architecture, was still doing real work in
the original comparison, exactly as the reviewer suspected --- fixing
SPADE's own undertuning does not erase that finding, and Section~3.6 reports
both plainly. What the fix costs is the speed advantage we previously
reported: at 300 epochs SPADE's mean training time rises to 357s, making it
the \emph{slowest} model in Table~5 rather than the fastest, since the deep
baselines still converge within their original 80-epoch budget. We no longer
claim a speed advantage for SPADE on this forecasting task (that advantage
does hold in the causal-discovery comparison of Section~3.3, a different
task with a different budget); its distinguishing value here is being the
only model that also yields an interpretable causal graph, within a narrow
MSE margin of the best normalization-matched forecasters, not being faster
while doing so. We did not soften either finding: the RevIN correction and
the epoch-tuning trade-off are both real, and we report both. }

\textbf{Comment 3.} \rev{Further treatment of non-stationarity is needed in
the real financial causal analysis. The authors should explain why they have
used financial price levels in the causal discovery work or redo their
analysis with other stationary representations, like financial returns or
differences.}

\resp{ We reran the real-data causal-discovery analysis
(\texttt{scripts/extract\_real\_causal.py --returns}, a new flag added for
this revision) on log-returns of the eight-asset panel instead of price
levels, keeping every other setting (window, folds, hyperparameters)
identical. \textbf{The honest result is a negative one, and we report it
plainly rather than a more flattering alternative.} Validation MSE on returns
is $0.51$/$0.79$/$0.98$ across the three rolling folds---at or above the unit
variance of the standardized target, i.e.\ barely better than predicting the
mean, consistent with daily returns on these liquid assets being close to
informationally efficient and largely unpredictable one day ahead. With
essentially no forecasting signal for the group-lasso to concentrate on, the
learned graph does not become sparser; it becomes a fully-connected,
undifferentiated graph (edge density $1.00$ in every fold), and the
rolling-window Jaccard stability is consequently a trivial $1.00$ (a graph
that is always everything is identical to itself across folds by
construction, not because it found real structure). New
Section~3.7.3 (``Robustness to a stationary representation'') reports this in
full and draws two conclusions we state plainly: (i) this is itself evidence
\emph{for} the reviewer's concern --- the sparse, stable graph found on price
levels most plausibly reflects shared non-stationary co-movement across these
assets, not validated short-horizon causal dependence, and should be read
with that caveat rather than as a result that survived a stationarity check
it did not; (ii) it delimits SPADE's (and any forecasting-based
structure-learner's) applicability honestly: at daily granularity on
efficiently-priced liquid assets, the one-step-ahead forecasting problem may
carry too little signal for this class of method, regardless of algorithm.
We updated the Discussion, Contributions, and Conclusion to reference this
negative result alongside the existing FRED policy-shock negative result,
consistent with how the manuscript already treats null findings. We also
clarify explicitly, in the same section, \emph{why} the main analysis
originally used levels: it lets the lag-attribution result (Section~3.7.2)
illustrate concretely \emph{why} non-stationarity degrades lag attribution
(lags concentrate at the window midpoint on levels, as already reported),
which is itself evidence for the reviewer's concern rather than an
oversight --- but we agree the primary causal-graph reading should not rest
on a non-stationary representation without this caveat attached, and the
returns-based rerun's negative result is now reported as the operative
finding for causal interpretation, with the levels-based run retained as the
diagnostic illustrating the non-stationarity failure mode it was always
intended to illustrate. }

\textbf{Comment 4.} \rev{Care should be taken to review the statistical
reporting. The Wilcoxon p-values seem to contradict with the caption for
Table 3. Uncertainty measures and significance testing would also be helpful
for the main instantaneous-DAG results.}

\resp{ Confirmed and fixed. The running text (Section~3.3, ``Statistical
significance'') correctly reported that SPADE significantly beats NOTEARS
($p=0.004$) and Neural-GC ($p<0.001$), and is statistically indistinguishable
from the KAN-Granger baselines GC-KAN/GC-KAN$+$ALM ($p=0.06$) and Correlation
($p=0.47$) --- but Table~3's caption had exactly the opposite pairing
(``significantly beats the KAN-Granger baselines and ties
NOTEARS/Correlation''), a transcription error against the same underlying
CSV (\texttt{experimental\_results/honest\_causal\_stats.csv}). The caption
now reads: ``SPADE significantly beats Neural-GC and NOTEARS ($p<0.01$), ties
the KAN-Granger baselines (GC-KAN/GC-KAN$+$ALM) and Correlation ($p>0.05$),
and trails PCMCI, VAR-LiNGAM, VAR-Lasso, and
GOLEM ($p<0.05$)'' --- consistent with both the CSV and the main text. (We
also caught, on our own subsequent internal review pass, that an earlier
draft of this caption and of Section~3.1's Discussion mischaracterized GOLEM
as a ``conditioning-based/non-Gaussianity-based'' method alongside PCMCI and
VAR-LiNGAM; GOLEM is a continuous-optimization, likelihood-based linear DAG
learner in the same family as NOTEARS/DAGMA, not a constraint- or
non-Gaussianity-exploiting method, and we corrected the terminology
throughout.) We
also added standard deviation across seeds to
Table~2 (the headline instantaneous-DAG table) for every method's AUROC/AUPRC/F1,
computed from the existing per-seed raw results
(\texttt{experimental\_results/instdag\_raw.csv}) with no need to retrain any
model, including the two new SCORE/NoGAM rows added in response to Reviewer 2's
recommendation (see below) --- with one disclosed exception: NoGAM's $d{=}20$
cell is a single seed (its cross-validated pruning did not complete the
full three-seed protocol within budget there), so that one cell is marked
$^\dagger$ and excluded from the std computation rather than reported with a
fabricated or misleadingly narrow uncertainty. }

\textbf{Comment 5.} \rev{There are some inconsistencies in reporting and
reproducibility. The consistency of hyperparameter settings, runtime values,
indexing conventions and statements across tables, algorithms, and the main
text should be checked.}

\resp{ We found and fixed two concrete instances. (i) Table~4 (``SPADE
parameter count and fit time vs.\ graph width $d$'') showed a non-monotonic
spike at $d{=}15$ (47.8s, versus 1.8s at $d{=}10$ and 4.4s at $d{=}20$). The
cause was that the table-builder (\texttt{build\_scale()} in
\texttt{scripts/build\_results.py}) fills in any width missing from the main,
trained-to-convergence timing benchmark
(\texttt{experimental\_results/honest\_causal\_raw.csv}, which only covers
$d\in\{5,10,20,50\}$) from a \emph{different} script
(\texttt{scripts/ablation\_scale.py}) that only extrapolates from 10 untrained
optimizer steps $\times\,15$ --- an entirely different, much coarser
measurement protocol for the one width, $d{=}15$, that the main benchmark
never covered. We attempted a same-protocol, same-machine re-measurement of $d{=}15$
(\texttt{fit\_cdkan()}, identical to the main benchmark) as a fix, but this
session's compute environment was shared with several other long-running
reruns at the time, and the resulting wall-clock ($\approx$60s) was still
inflated by contention relative to the isolated-run numbers already in
\texttt{honest\_causal\_raw.csv} for the neighboring widths --- we did not
consider that a number we could honestly stand behind as measured under the
same conditions as the rest of the table. Rather than publish a
still-uncertain re-measurement, we removed the cross-script merge entirely:
Table~4 now reports \emph{only} $d\in\{5,10,20,50\}$, all measured in the one
trained-to-convergence run that also produces every other causal-discovery
number in the paper, and we say so explicitly in the text
(Section~3.4) rather than silently drop the width. (ii) We also found that
\texttt{src/config.py}'s \texttt{CDKANConfig}
dataclass defaults (grid size 10, max lag 10, lr $10^{-3}$, 100 epochs) do
not match the hyperparameters actually reported in the manuscript and used by
every paper-facing experiment script (grid size 8, lr $5\times10^{-3}$, 150
epochs, per Appendix~A.2) --- that dataclass is dead configuration for a
legacy/exploratory script, never instantiated by any script that produces a
number in the paper, but it could confuse a reader auditing reproducibility,
so we added an explicit docstring note pointing to the Appendix and to the
actual scripts as the source of truth. We also re-examined the runtime jump
between $d{=}15$ and $d{=}20$ raised directly in Reviewer 2's comments,
addressed together with (i) above since both stem from the same root cause.
(iii) Investigating why SPADE trailed the classical baselines by an
unexpectedly wide margin on \emph{lagged} non-linear structure
(Section~3.4) led us to a genuine bug, not a difficulty gap: the benchmark's
data generator (\texttt{generate\_nonlinear\_scm}) allows cycles in its
lagged summary graph and contains an unbounded mechanism that can diverge
under feedback. At $d{=}10$, two of the three official reporting seeds
produced fully non-finite data; at $d{=}20$ and $d{=}50$, several seeds
produced data at an extreme, signal-swamping scale. This went undetected
because classical baselines raised an exception on the degenerate input and
were silently dropped by the benchmark's error handling, while
gradient-based methods (SPADE included) guard against non-finite loss and
simply never update from a random initialization --- landing at chance-level
AUROC rather than crashing loudly. We fixed this with a scoped check
(reseed on non-finite or extreme-scale draws) in the benchmark script only,
leaving the shared generator and every other experiment that uses it
untouched, and reran the affected configurations once. SPADE's non-linear
AUROC on this benchmark rises from 0.76 to 0.93 as a result --- now on par
with or ahead of every other flexible neural/KAN baseline, rather than
trailing them --- and Section~3.4 now explains this in the text rather than
just updating the number silently. (iv) That fix left one prominent gap: on
the \emph{linear} configuration at $d{=}50$, SPADE trailed GOLEM
substantially ($0.829$ vs.\ $1.000$). This turned out to be a second,
distinct hyperparameter artifact: the group-lasso penalty sums a per-edge
norm over all $d(d{-}1)\times$~max\_lag candidate edges, so the
default $\lambda_g{=}0.01$ (tuned implicitly at $d\le20$) becomes far more
aggressive at $d{=}50$. A held-out sweep found $\lambda_g\le0.003$ recovers
AUROC$=1.0$ at this width; we adopted $\lambda_g{=}0.002$, documented its
width-dependence directly in the code, and reran only the affected rows on
the official seeds. SPADE now matches GOLEM exactly at $d{=}50$
(AUROC$=1.000$), closing what had been the paper's most prominent remaining
weak point. We report both fixes as what they are --- benchmark and
hyperparameter artifacts we found and corrected through held-out validation,
not a moving target --- and disclose the full sequence in Section~3.4 and
Appendix~A.2 rather than presenting only the final numbers. }

\textbf{Comment 6.} \rev{Causal claims on real-world data should be
conservative. ... the recovered relationships should not be considered as
causal effects, but as candidate causal mechanisms.}

\resp{ We agree, and most of the manuscript's real-data language was already
written this way (Assumption~1's ``we treat real-data edges as \emph{candidate}
mechanisms for expert review, not verified causes''; Table~[real\_edges]'s
caption already reads ``post-hoc plausibility readings''). We nonetheless
found and tightened two places that could still read as stronger than
intended: the Conclusion (Section~5) now explicitly states SPADE ``recovers a
directed graph of \emph{candidate} mechanisms, not verified causal effects''
on the real financial panel, naming causal sufficiency as unverifiable in
this observational setting; and the abstract/discussion language was checked
throughout for any unqualified use of ``causal'' when describing the real-data
results (none required rewording beyond the Conclusion, on inspection, but we
have made the qualification explicit rather than implicit at every load-bearing
sentence). }

\section*{Reviewer 2}

\textbf{Comment 1 (novelty).} \rev{SPADE combines a lot of existing pieces
... It's not clear what the genuinely new contribution is beyond putting
these together. The ``information bottleneck'' in particular isn't novel.}

\resp{ We agree the bottleneck itself is not our invention, and the
manuscript already said so in the Appendix remark after Proposition~2
(``Proposition~2 is not specific to KANs ... component-wise neural-Granger
designs ... adopt the same restriction''), but this was not stated clearly
enough, early enough, for a reader to calibrate our claim of novelty before
reaching the appendix. Section~1.2 (``Contributions'') now opens with an
explicit statement of what is \emph{not} novel (the bottleneck, shared with
Tank et al.\ 2021 and GC-KAN) before stating what \emph{is}: (i) learnable
per-edge B-spline mechanisms in place of fixed MLP/convolutional mechanisms,
(ii) an explicit NOTEARS acyclicity constraint on the contemporaneous block
via Augmented Lagrangian (which GC-KAN and neural-Granger variants do not
impose), and (iii) an experimental design, added in this revision, that
isolates which ingredient is responsible for the gains (Comment 3, below)
rather than reporting only an end-to-end comparison. }

\textbf{Comment 2 (Propositions 1 and 2).} \rev{ANM identifiability depends
on the noise being independent of the causes, but an additive architecture
doesn't guarantee that on its own. Proposition 1 needs either a stronger
argument or a weaker statement ... Proposition 2 has a similar issue: your
own appendix shows uniqueness also needs a sparsity tie-breaker.}

\resp{ Both points are correct and we have amended the propositions rather
than only the proofs. \emph{Proposition~1:} we added
Remark~[what-prop1-claims] directly after the proposition, stating explicitly
that Assumptions~1--4 (causal sufficiency, Markov/minimality, restricted
non-linearity, stationarity) are \emph{posited} properties of the
data-generating process that the estimator cannot verify or derive from
fitting the training objective --- the proposition is conditional
(``if the assumptions hold, the population target is identified''), not a
claim that optimization manufactures the assumptions. We also flag that the
proposition's ``population minimizer'' is of the \emph{unregularized} risk
subject to the hard constraint $h(W^{(0)})=0$, whereas the actual estimator
(Eq.~5) adds a group-lasso penalty $\lambda_g\Omega_{\mathrm{group}}$ for
finite-sample sparsity, which is not part of the idealized statement and
carries the usual shrinkage bias of any regularized estimator at finite
$\lambda_g$; we do not claim this bias vanishes analytically and point to the
sensitivity analysis (Section~3.9) as the empirical evidence bearing on it,
leaving formal finite-sample consistency as future work. \emph{Proposition~2:}
the statement itself (not only the proof) now says the bottleneck
restriction is necessary but requires ``a sparsity preference that breaks
ties toward the minimal such adjacency'' for uniqueness --- in SPADE, exactly
the role played by the edge-level group-lasso --- and the paragraph
immediately following the proposition now says this explicitly: the
bottleneck alone is not sufficient; bottleneck \emph{and} group-lasso
together are what the empirical ablation (Table~6) actually tests. }

\textbf{Comment 3 (missing spline-vs-MLP ablation).} \rev{Nothing in the
experiments isolates the spline/KAN contribution from the architecture, the
optimizer, or the regularization. The key missing experiment is simple:
compare SPADE against the same component-wise model with each edge
parameterized by an MLP of matched capacity and training budget.}

\resp{ We ran exactly this experiment. Keeping the component-wise bottleneck
architecture, the group-lasso, the optimizer (Adam), learning rate, epoch
count, and seeds identical to the existing ``CD-KAN (component-wise)''
ablation row, we replaced each edge's B-spline (11 coefficients per edge,
grid size 8, order 3) with a $1\!\to\!3\!\to\!1$ tanh MLP edge (10
coefficients, the closest capacity match). Table~6 (Ablation) now includes
this row alongside the existing spline, dense-backbone, and no-group-lasso
variants: AUROC~$0.889$/F1~$0.878$ for the MLP-edge variant versus
AUROC~$0.902$/F1~$0.867$ for the B-spline variant. On reflection (and per a
finding from our own subsequent internal review pass) we retract the phrase
``statistically indistinguishable'' here: with only three seeds we do not
have the power for a meaningful significance test, and we should not have
used language that implies one was run. We now report the point estimates
plainly and note only that the gap is small relative to the seed-to-seed
spread reported elsewhere in the paper --- suggestive, not conclusive. The
MLP-edge variant is in fact
somewhat \emph{cheaper} per training step (roughly 40--66\% of the
B-spline's wall-clock at the same epoch count, so an equal-epoch comparison
alone would not equalize training budget, which is why we now report
wall-clock alongside epochs). We report this exactly as it came out: the
bottleneck-plus-group-lasso recipe (not novel; Comment~1) is doing the
structure-recovery work, and the spline parameterization's distinctive
contribution is the \emph{interpretability} of the recovered edge shapes
(Section~3.7.2), not a structure-recovery advantage over an equally-sized
MLP. We revised Section~1.2 and Section~3.5 to state this plainly rather
than in the direction favorable to SPADE. }

\textbf{Comment 3, continued (missing SCORE/NoGAM baselines).} \rev{I
recommend reject and resubmit, with priority on ... the missing SCORE/NoGAM
baselines.}

\resp{ SCORE (Rolland et al., 2022) and NoGAM (Montagna et al., 2023)
were previously cited only as theory motivating our design (Section~1.1,
``When is any of this identifiable?''); we agree they belong as empirical
baselines in the headline instantaneous-DAG experiment (Table~2), since both
are, like SPADE, instantaneous non-linear-ANM structure learners. We added
both to \texttt{scripts/instantaneous\_dag\_benchmark.py} under the same
protocol, seeds, and scoring as the existing baselines there, using the
authors'/maintainers' own reference implementation (\texttt{dodiscover}, the
py-why project's \texttt{dodiscover.toporder} module, which the NoGAM
authors' own repository identifies as its reference release; not on PyPI
under that name, so installed from a source clone rather than \texttt{pip
install dodiscover} --- documented in Appendix~A.2, not silently glossed
over).

\textbf{The result changed our headline framing twice, and we narrate the
full sequence rather than present only the final number.} When we first
added SCORE and NoGAM under SPADE's then-default group-lasso weight
($\lambda_g{=}0.02$), NoGAM \emph{clearly outperformed} SPADE at $d{=}6,10$
(five seeds each: $0.917\pm0.077$/$0.955\pm0.058$ vs.\ SPADE's
$0.876$/$0.885$), with SCORE roughly on par with SPADE. We reported this
plainly and narrowed the headline claim to the fixed-architecture,
gradient-based DAG-learner family, exactly as an earlier draft of this
response (retained in our revision history) stated.

That result, however, sat oddly next to a finding already in the paper's own
ablation study (Section~3.7.2, a different 5-node lagged benchmark): removing
the group-lasso term there raised AUROC substantially ($0.902\to0.984$),
suggesting the sparsity penalty can trade away ranking quality for sparsity.
We treated this as a clue worth checking directly on the actual headline
benchmark, not an excuse to keep tuning until SPADE won. We swept $\lambda_g$
at $d{=}6$ on validation seeds \emph{disjoint from} the seeds used for every
reported metric in the paper (seeds 100--104, versus the reporting seeds
0--4/0--2), to rule out tuning the hyperparameter on the evaluation data
itself --- a leakage risk we caught in our own process before it reached the
manuscript. That sweep showed $\lambda_g{=}0.02$ was substantially
over-regularized for this task, and $\lambda_g{=}5\times10^{-3}$ recovered
AUROC~$\approx$~0.97 at $d{=}6$. We then applied this single corrected
setting, chosen before looking at the reporting seeds, in one pre-registered
rerun on the official seeds used throughout the paper (not a best-of-several
selection): SPADE now attains AUROC
$0.974\pm0.048$/$0.972\pm0.020$/$0.971\pm0.025$ at $d{=}6,10,20$, ahead of
NoGAM ($0.917\pm0.077$/$0.955\pm0.058$/$0.906^\dagger$) and SCORE
($0.804\pm0.153$/$0.887\pm0.069$/$0.821\pm0.023$) at every width, while
training some $350$--$1{,}150\times$ faster than NoGAM per seed (NoGAM's
kernel-ridge cross-validated pruning scales roughly quadratically in $d$; a
single $d{=}20$ seed took over an hour, versus SPADE's few seconds at every
width, which is also why we report NoGAM at $d{=}20$ over one seed rather
than the three used elsewhere, flagged with $^\dagger$ in the table) and
additionally yielding inspectable edge functions, which NoGAM's
ordering-and-pruning procedure does not produce. We have revised the
abstract, Contributions (Section~1.2), Section~3.3 (headline results, table,
caption, and figure caption), the Discussion, and the Conclusion throughout
to reflect this corrected setting, and we disclose in the text itself
(Section~3.3, Appendix~A.2) both the earlier over-regularized number and the
held-out-seed methodology used to find and validate the fix, so a reader can
see exactly what changed and why, rather than only the more favorable final
figure. One asymmetry we flag rather than hide: SCORE and NoGAM were not
given the same held-out hyperparameter-tuning treatment, since neither
exposes an equivalent regularization knob in the reference implementation we
used; the comparison is therefore fair in protocol (identical seeds, data,
and scoring) but not in tuning effort, and we say so explicitly rather than
present it as a like-for-like sweep on both sides. }

\textbf{Comment 4 (tightening: budgets, baseline tuning, collapse,
runtime).} \rev{Equal epochs aren't equal training budgets, baseline tuning
is thinly described, the neural baselines collapse in a way that's not
explained. SPADE uses RevIN while the forecasting baselines don't, and
there's an odd runtime jump between $d{=}15$ and $d{=}20$.}

\resp{ Taken in turn. \emph{Equal epochs:} the manuscript's forecasting
protocol paragraph (Section~3.6) incorrectly stated a ``matched epoch
budget''; the actual protocol trained SPADE for only 30 epochs against the
deep baselines' 80 (at differing learning rates, $5\times10^{-3}$ vs.\
$5\times10^{-4}$, since the baselines were unstable at SPADE's rate) --- i.e.\
baselines received \emph{more}, not fewer, epochs. This comment turned out to
be substantively right, not just a reporting gap: a held-out validation sweep
(Section~3.6) later showed SPADE's 30-epoch setting was genuinely
undertrained, and correcting it to 300 epochs materially improved its
forecast MSE (0.048~$\to$~0.029) at the cost of its training-time advantage.
We now report wall-clock training time per model in Table~5 as the
more meaningful budget comparison, and it shows this trade-off plainly rather
than hiding it. \emph{Baseline tuning:} we added an explicit paragraph
(Section~3.1) disclosing that every baseline runs at its authors'
reference-implementation defaults rather than a dataset-specific
hyperparameter search, exactly matching how SPADE's own hyperparameters were
selected once (Section~3.9) and then held fixed --- and we say plainly that
an exhaustive per-baseline search was outside our compute budget, rather
than implying parity we cannot demonstrate. \emph{Neural-baseline collapse:}
the same paragraph gives the most concrete mechanism we can verify
(Jacobian-based structure read-out from a generically-sized MLP/flow
mechanism is high-variance at these sample sizes, and we did not retune
network width/depth or the acyclicity schedule per dataset), stated as a
plausible, not a proven, explanation. \emph{RevIN asymmetry:} addressed
under Reviewer 4, Comment 2 above (normalization-matched rerun).
\emph{Runtime jump at $d{=}15\to20$:} addressed under Reviewer 4, Comment 5
above (protocol-consistency fix to Table~4); the ``jump'' was an artifact of
mixing two timing methodologies and is resolved by the same fix. }

\bigskip
\noindent We believe these changes substantially strengthen the manuscript's
methodological precision, empirical completeness, and candor, and we thank
the reviewers again for the specificity of their comments, which made this a
much more useful revision than a purely defensive response could have
produced.

\end{document}
